
\section{Introduction}


\begin{frame}
\frametitle{Le partitionnement (\textit{sharding})}

Le partitionnement consiste à  

\begin{itemize}
  \item \alert{découper} une (grande) collection
      en \alert{fragments} en fonction d'une \alert{clé}\,;
      \item placer chaque fragment sur un \alert{serveur}\,;
      \item Maintenir un \alert{répertoire} indiquant  que telle \alert{clé} se trouve dans tel \alert{fragment}
          sur tel \alert{serveur}
\end{itemize}
 
\vfill

Le partitionnement apporte la \alert{répartition de charge} (\textit{load balancing})

\vfill

Il doit être \alert{dynamique} (ajout/retrait de serveurs) pour s'adapter ``élastiquement'' 
 
 \vfill
 
 S'applique à des collections de paires (clé, valeur), où valeur est n'importe quelle information structurée.
\end{frame}

\begin{frame}
\frametitle{Architecture d'un système avec partitionnement}

\figSlide{archi-partition}
      
\begin{itemize}
  \item Routeur = \textit{Master}, \textit{Balancer}, \textit{Config server}, ...
  \item Fragment = \textit{chunk}, \textit{tablet}, \textit{region}, \textit{bucket}, ...
\end{itemize}

\begin{remark}
Chaque serveur gère une réplication locale pour tolérance aux pannes.
\end{remark}
\end{frame}


\begin{frame}
\frametitle{Techniques de partitionnement}

\begin{itemize}
   \item par \alert{intervalle}\,: documents triés sur la clé, puis groupés par intervalles.
   
     Exemples représentatifs\,: \mongodb/, HBase (BigTable).
     
   \item par \alert{hachage} sur la clé, avec répertoire de hachage.
   
       Exemples\,: \textit{memcached}, \textit{chord}, \textit{Dynamo/S3/Voldemort}, beaucoup d'autres.
 \end{itemize}
 

\begin{remark}
Se référer aux techniques d'indexation par arbre-B et par hachage dans des environnements centralisés.
\end{remark}
\end{frame}

\begin{frame}
\frametitle{Opérations}

Les opérations de type ``dictionnaire'' peuvent être transmises à \alert{un seul} serveur\,:

\begin{itemize}
   \item $get(k): v$, recherche par clé
   \item $put (k, v)$, insertion
   \item $delete(k)$, suppression
   \item $range(k_1, k_2)$, recherche par intervalle (peut impliquer plusieurs serveurs, ne marche pas avec le hachage)
\end{itemize}

\vfill

\alert{Pour toutes les autres opérations}\,: tous les serveurs effectuent l'opération localement.

\vfill

Exemple\,: recherche sur une clé secondaire, s'effectue sur \alert{tous} les serveurs, avec si possible des index locaux.
\end{frame}

\begin{frame}
\frametitle{Notre application}

On veut gérer (légalement) une collection de vidéos.

\begin{itemize}
  \item chaque vidéo occupe en moyenne 500 MO
  \item on associe chaque vidéo à quelques méta données (titre, auteur, année, etc.)
  \item 1\,000 vidéos = 500 GO (1/2 TO), pas de souci avec un (bon) serveur.
  \item 1\,000\,000 vidéos = 500 TO (1/2 PO), il faut distribuer.
\end{itemize}

\begin{remark}
Dès qu'on parle de système distribué, il faut envisager le coût en stockage de la réplication. Multiplier
par 2 ou 3...
\end{remark}


\begin{noteblock}{Besoins}
Chercher une vidéo par son titre, les vidéos par année, intervalle d'années, etc.
\end{noteblock}

\end{frame}


\begin{frame}
\frametitle{Pour illustrer}

Un tout petit échantillon.

\begin{center}
\begin{small}
\begin{tabular}{r|l|c|c||r|l|c|c}
id &titre & année  & mp4  & id & titre & année  & mp4  \\
\hline
1 & Vertigo  & 1958 & ... & 9 & Annie Hall & 1977 & ... \\
2 & Brazil & 1984 & ... & 10&Jurasic Park & 1992 & ... \\
3&Twin Peaks & 1990 & ...&  11&Metropolis & 1926 & ... \\
4&Underground & 1995 & ... & 12&Manhattan & 1979 & ... \\
5&Easy Rider & 1969 & ... & 13&Reservoir Dogs & 1992 & ... \\
6&Psychose & 1960 & ... & 14&Impitoyable & 1992 & ... \\
7&Greystoke & 1984 & ... & 15&Casablanca & 1942 & ... \\
8&Shining & 1980 & ... & 16&Smoke & 1995 & ... \\
\hline
\end{tabular}
\end{small}
\end{center}

Quelle clé considère-t-on\,? On va y revenir.

\end{frame}
\section{Partitionnement par intervalle}

\subsection{Quelques rappels}


\begin{frame}
\frametitle{Première étape\,: il faut trier}

Il faut disposer d'une structure triant les documents sur la clé. Deux possibilités\,:

\vfill

\alert{1) On trie la collection (clé, document)}

\begin{itemize}
  \item indexation dite ``plaçante'': l'emplacement d'un document est déterminé par sa clé.
  \item efficace en recherche (pas d'indirection)
  \item peu efficace en mise à jour (nécessite des déplacements)  
  \item on peut créer un seul index plaçant par collection
\end{itemize}

\alert{2) On trie une structure indépendante constituée des paires (clé, adresse) }

\begin{itemize}
  \item indexation non-plaçante\,: les documents peuvent être insérés n'importe où (efficace)
  \item un peu moins efficace en recherche (on trouve une adresse dans l'index, pas le document lui-même)
  \item autant d'index que l'on veut
\end{itemize}

\end{frame}



\begin{frame}
\frametitle{Seconde étape (optionelle)\,: on crée des niveaux d'index}

Illustration avec \alert{collection triée}, la clé est le titre.

\figSlideWithSize{idxnondense}{10cm}

On a découpé en \alert{4 fragments}, 
indexés avec un niveau.

\end{frame}


\begin{frame}
\frametitle{Quelques calculs, cas d'une collection triée}

On a un million de vidéos, soit $2^{20}$, \alert{triés} sur la clé.

\vfill

\alert{Sans index}
\begin{itemize}
  \item Recherche par dichotomie, soit 20 accès disques pour trouver une vidéo
  \item Temps\,: 20 fois 10 ms = 200 ms (au pire du pire)
\end{itemize} 

\alert{Avec index}
\begin{itemize}
  \item Mon index comprend 1 million de paires (clé, adresse)
  \item Supposons qu'une adresse occupe 8 octets, la clé 8 octets aussi\,: l'index occupe 16 MO 
  \item Temps\,: recherche quasi instantanée dans l'index, 1 accès au disque (10 ms au pire) 
\end{itemize} 

\begin{noteblock}{Conclusion}
L'index est tout petit (tient en mémoire), permet des accès très rapides même sur de très grandes collections.
\end{noteblock}
\end{frame}


\begin{frame}
\frametitle{A quoi sert de \alert{trier} la collection\,?}

Différence (subtile)\,: si la collection est triée, je peux créer
un \alert{index non-dense}, sinon l'index \alert{doit} être dense.

\vfill
\alert{Index non-dense}

\begin{itemize}
  \item Je découpe ma collection en \alert{fragments} (quelques MO ou dizaines de MO)
  \item Chaque fragment couvre un intervalle
  \item \alert{J'indexe seulement les intervalles}
\end{itemize} 

Exemple dans le cas précédent: je découpe ma collection en fragments de 4 films; j'indexe
250\,000 fragments (1M/4).

\vfill

\alert{Si la collection n'est pas triée, je dois créer un index dense} (pas le choix)
\begin{itemize}
  \item Chaque document est représenté dans l'index
\end{itemize} 

\end{frame}

\begin{frame}
\frametitle{Exemple d'un index dense}

On indexe sur l'année, en y mettant \alert{toutes} les clés.

\figSlideWithSize{idxdense2}{10cm}

Presque aussi efficace, sauf pour les recherches par intervalle.
\end{frame}

\subsection{Indexation distribuée}


\begin{frame}
\frametitle{Appliquons ces principes à la distribution des données}

A quelques variantes près, on procède comme pour un index non-dense.

\begin{itemize}
  \item La collection est triée sur la clé.
   
        \smallskip
   
       Au moins partiellement (cf. plus loin)
       
  \item On la découpe en fragments d'une taille maximale pré-déterminée (de 64MO à qq GO)
   
   \smallskip
   
        Chaque frament couvre donc un intervalle $[min_f, max_f[$
        
   \item Les fragments sont distribués équitablement sur les serveurs

    \smallskip
    
         Un ou plusieurs fragments par serveur.

   \item Un index sur les fragments est maintenu par le routeur.
 
    \smallskip
   
      \alert{Souvenez-vous\,: l'index est tout petit par rapport à la collection}
      
\end{itemize} 

\alert{Le routeur reçoit les requêtes et détermine le(s) serveur(s) à solliciter.}

\end{frame}

\begin{frame}
\frametitle{L'architecture, avec partionnement par intervalle}

Le routeur connaît les intervalles, et les serveurs associés.

\figSlide{partition-range}

\begin{remark}
\begin{itemize}
   \item Chaque serveur maintient un index local sur ses fragments.
   \item Les fragments sont autant que possible conservés en mémoire.
  \item Parmi les variantes\,: les fragments sur les serveurs peuvent être triés ou pas. 
\end{itemize}
\end{remark}
\end{frame}


\begin{frame}
\frametitle{Dynamicité: comment évolue le système\,?}

L'opération de base est le \alert{split}, division d'un gros fragment en deux petits.

\figSlide{partition-split}

\alert{L'équilibrage} consiste à répartir équitablement les fragments.

\end{frame}

\begin{frame}
\frametitle{Les opérations du routeur}

% A two-columns presentation
\begin{tabular}{l|p{6cm}}
   \parbox{5cm}{
$get(k)$: chercher l'intervalle qui contient $k$, transmettre
au serveur correspondant.

\smallskip

$put(k, d)$: identifier le serveur avec $k$, transmettre l'insertion
 
 \smallskip
 
 $delete(k)$: idem
 
 \smallskip
 
 $range(k_1, k_2)$: transmettre à tous les serveurs gérant un intervalle chevauchant $[k_1, k_2]$
 
 \smallskip
 
 Toute autre opération\,: transmettre à tous les serveurs.
   }  % Inclusion of an XML file
  & 
  \parbox{6cm}{
 
\begin{tabular}{r|l}
Intervalle  & Serveur\\
\hline
$]-\infty, a]$ & A  \\
$[a+1, b]$ & B  \\
$[b+1, c]$ & C  \\
$[c+1, d]$ & B  \\
$[d+1, \infty[$ & A  \\
\hline
\end{tabular}

}

\end{tabular}

\end{frame}


\section{Etude de cas 1\,: BigTable (HBase, Cassandra, ...)}

\begin{frame}
\frametitle{BigTable\,: la structure} 

C'est une collection \alert{triée}, dont chaque ligne est indexée par une clé.


 \figSlideWithSize{bigtable}{8cm}
 
La table est partionnée en fragments (\textit{tablets}) qui sont eux-mêmes indexés 
par une racine. 

\vfill

Les fragments pleins subissent un \alert{split}, le nouveau fragment est propagé 
au père.


\end{frame}
 

\begin{frame}
\frametitle{Le routeur}

Le routeur est incorporé au \textit{driver} et maintient une \alert{image} de l'arbre.


\figSlideWithSize{clientImage}{7cm}

Cette image peut être \alert{partielle} et \alert{obsolète}. Elle est \alert{ajustée}
à chaque requête.

\end{frame}
 
\begin{frame}
\frametitle{Exemple d'un ajustement}

Une requête Client transmise sur la base de l'image échoue, car l'arbre a 
beaucoup évolué.

\figSlide{imageadjust}

L'ajustement consiste à remonter/descendre dans l'arbre.

\end{frame}

\section{Etude de cas 2\,: \mongodb/}

 
\begin{frame}
\frametitle{Le \textit{sharding} de \mongodb/}

On crée des \textit{replica set} qui vont chacun stocker un sous-ensemble des fragments.

\vfill

Un serveur spécial, le \textit{config server} (répliqué), se charge de stocker la table de routage.

\vfill

Chaque application discute avec un processus routeur, \textit{mongos} 

\vfill

\figSlideWithSize{mongosharding}{8cm}


\end{frame}



\begin{frame}[fragile]
\frametitle{Allons-y pour un petit test}

On lance un \textit{config server}

\begin{smallverbatim}
mkdir /data/configdb
mongod --configsvr --dbpath /data/configdb --port 27019
\end{smallverbatim}

On lance un routeur en lui indiquant où se trouve(nt) le(s) \textit{config server}(s)

\begin{smallverbatim}
mongos --configdb localhost:27019
\end{smallverbatim}

Et deux serveurs (il faudrait des \textit{replica sets} complets en production)
\begin{smallverbatim}
mongod --dbpath /data/node1 --port=30001
mongod --dbpath /data/node3 --port=30002
\end{smallverbatim}

A partir de \textit{mongos}, déclarons les deux serveurs comme \textit{shards}.

\begin{smallverbatim}
mongo --port 27017
sh.addShard("localhost:30000")
sh.addShard("localhost:30001")
sh.enableSharding("nfe204")
\end{smallverbatim}

\end{frame}

\begin{frame}[fragile]
\frametitle{Choisir la clé de partitionnement}

On indique la clé de partitionnement avec la commande suivante\,:

\begin{smallverbatim}
sh.shardCollection("nfe204.videos", { "title": 1 } )
\end{smallverbatim}

Attention au choix de la clé

\begin{itemize}
  \item Les clés-séquences\,:  on envoie toujours les insertions au même serveur\,; bon ou pas?
  \item Les clés à faible cardinalité (ex: pays)\,: pas très bon non plus, car 
               il pourra être difficile de faire un \textit{split}
\end{itemize}

Bonne option: appliquer un hachage pour bien distribuer les insertions\,:

\begin{smallverbatim}
sh.shardCollection("nfe204.videos", { "title": "hashed" } )
\end{smallverbatim}

\end{frame}


\begin{frame}[fragile]
\frametitle{Inspecter la configuration}

Liste des \textit{shards}
\begin{smallverbatim}
db.runCommand({listshards: 1})
\end{smallverbatim}

Statut général du cluster
\begin{smallverbatim}
sh.status()
\end{smallverbatim}

Les fragments sont dans une collection \textit{chuncks}.

\begin{smallverbatim}
use config
db.chuncks.count()
db.chuncks.findOne()
\end{smallverbatim}
\end{frame}

%%%%%%%%%%%%%%%%%%%%%%%%%%%%%%%%%%%%%%%
\section{Partitionnement par hachage}
%%%%%%%%%%%%%%%%%%%%%%%%%%%%%%%%%%%%%

\subsection{Rappels}


\begin{frame}
\frametitle{A la base\,: table de hachage et fichiers}

On \textit{calcule} la position d'un enregistrement
d'après la clé.
\begin{itemize}
  \item un \textit{fonction de hachage} $h$ associe des
         valeurs de clé à des adresses de bloc
   \item $h$ doit répartir uniformément les enregistrements
            dans les $n$ blocs alloués à la structure
  \item recherche d'un enregistrement =$>$ calcul de l'endroit
          où il se trouve.
\end{itemize}
Simple et efficace\,!
\end{frame}


%------------------------------------------------------------
\begin{frame}{Exemple}

On veut créer une structure de hachage pour nos 16 films
(hyp.\,: 4 enregistrements   par page)
\begin{itemize}
  \item on alloue 5 pages (pour garder une marge de man{\oe}uvre)
  \item un répertoire à 5 entrées (0 à 4) pointe vers les pages
   \item On définit la fonction $h(titre) = rang(titre[0])\, mod\, 5$
\end{itemize}

\end{frame}

%------------------------------------------------------------
\begin{frame}{Le résultat}

  \figSlideWithSize{hachage}{10cm}

\end{frame}


\begin{frame}{Recherches}

\begin{itemize}
  \item Par clé\,? Oui:\\
     \texttt{SELECT * FROM  Film WHERE titre = 'Impitoyable'}

   \item Par préfixe\,? Ici, oui.\\
      \texttt{SELECT * FROM  Film WHERE titre LIKE 'M\%'}

   \item Par intervalle\,: non\,!\\

\texttt{SELECT *
FROM  Film
WHERE titre BETWEEN 'Annie Hall' AND 'Easy Rider'}
  
\end{itemize}

\end{frame}
%------------------------------------------------------------
\begin{frame}{Mises à jour\,: ça se gâte}

La structure simple décrite précédemment n'est pas
\textbf{dynamique}
\begin{itemize}
  \item On ne peut pas changer un enregistrement de place

   \item Donc  il faut créer un chaînage de pages quand
             une page déborde

  \item Et donc les performances se dégradent...
\end{itemize}

\end{frame}
%------------------------------------------------------------
\begin{frame}{Exemple\,: insertion de Citizen Kane}

  \figSlideWithSize{hachage-2}{10cm}

\end{frame}

%------------------------------------------------------------
\begin{frame}{Hachage dynamique}

Objectif\,: réorganiser la table de hachage en fonction
des insertions et suppressions. 

\begin{itemize}
  \item le nombre d'entrées dans le répertoire est une puissance
          de 2
  \item la fonction $h$ donne toujours un entier sur 4 octets (32 bits)  
\end{itemize}
Idée de base\,: 
 on utilise les $n$ premiers bits du résultat de la fonction, avec
$n < 32$

\end{frame}
%------------------------------------------------------------
\begin{frame}{Exemple\,: le hachage des 16 films}

\begin{center}
\begin{tabular}{l|c}
titre & $h$(titre)   \\
\hline
Vertigo     & 01110010  \\
Brazil      & 10100101  \\
Twin Peaks  & 11001011  \\
Underground & 01001001  \\
Easy Rider  & 00100110 \\
Psychose    &  01110011 \\
Greystoke   &  10111001 \\
Shining     &  11010011 \\
\hline
\end{tabular}
\end{center}

\end{frame}
%------------------------------------------------------------
\begin{frame}{Construction de la table}

On départ on utilise seulement le premier bit
de la fonction
\begin{itemize}
  \item Deux valeurs possibles\,: 0 et 1
  \item Donc deux entrées, et deux blocs
  \item L'affectation d'un enregistrement dépend
        du premier bit de sa fonction de hachage  
\end{itemize}

=$>$ pour l'instant on reste dans un cadre classique

\end{frame}
%------------------------------------------------------------
\begin{frame}{Avec 5 films}

  \figSlide{hashext-1}

\end{frame}
%------------------------------------------------------------
\begin{frame}{Insertions}

Supposons 3 films par bloc. L'insertion de Psychose
(valeur 01110011) entraîne le débordement du premier bloc.

\begin{itemize}
  \item On double la taille du répertoire
  \item On alloue un nouveau bloc pour l'entrée 01
  \item Les entrées 10 et 11 pointent sur le \textit{même} bloc.
\end{itemize}
On agrandit seulement ce qui est nécessaire.

\end{frame}
%------------------------------------------------------------
\begin{frame}{Illustration}

  \figSlide{hashext-2}

\end{frame}
%------------------------------------------------------------
\begin{frame}{Insertions suivantes}

Plusieurs cas
\begin{itemize}
  \item on insère dans un bloc plein, mais plusieurs
           entrées pointent dessus
  \begin{itemize}
    \item on alloue un nouveau bloc, et on répartit les pointeurs
  \end{itemize}
  \item on insère dans un bloc plein, associé à une entrée
  \begin{itemize}
    \item on double à nouveau le nombre d'entrées
  \end{itemize}
Problème\,: le répertoire peut devenir très grand.
\end{itemize}

\end{frame}
%------------------------------------------------------------
\begin{frame}{Greystoke 10111001 et Shining 11010011}

  \figSlide{hashext-3}

\end{frame}



%%%%%%%%%%%%%%%%%%%%%%%%%%%%%%%%
\subsection{Hachage distribué}


\begin{frame}
\frametitle{Systèmes distribués basés sur le hachage}

Idée simple\,: tout le monde utilise la même fonction de hachage, les ``blocs''
sont les serveurs, et tout marche bien.

\vfill

\alert{Presque\,!} A ceci près que\,:

\begin{itemize}
  \item le système est nécessairement \alert{dynamique}\,; il faudra ajouter ou retirer
  des serveurs.

\item attention aux clients qui pouraient utiliser une fonction de hachage \alert{obsolète}\,: il faut
   leur garantir que leurs requêtes aboutissent quand même.
   
\end{itemize}

\vfill

En 1997, un article a proposé une solution maintenant très largement utilisé\,: le
\alert{hachage cohérent} ou \textit{Consistent haching}.


\begin{itemize}
  \item D'abord conçu et utilisé pour des \alert{caches distribués} (\textit{memcached})

  \item Popularisé pour la gestion de données par le système Dynamo (Amazon, 2007)

  \item Maintenant intégré à de très nombreux systèmes\,: Voldemort, Riak, Chord, ...
 
\end{itemize}
\end{frame}

\begin{frame}
\frametitle{Consistent hashing}

Soit $N$ le nombre de serveurs. La fonction suivante 
$$hash(key) \to modulo (key, N) = i$$
associe une paire $(key, value)$ au $S_i$.

\vfill

\alert{Problème}: si $N$ change, ou si un client utilise un valeur incorrecte de $N$, le hachage devient
\alert{incohérent}.

\vfill

Avec le  \alert{Consistent hashing}, ajouter ou retirer un serveur n'affecte la collection
que de façon \alert{marginale}.

\begin{itemize}
  \item on prend une fonction \alert{permanente}, $h$, appliquée aux clés \alert{et} aux serveurs
          vers l'espace d'adressage  $A=[0, 2^{64}-1$];
  \item $A$ est organisé comme un anneau parcouru dans le sens des aiguilles d'une montre\,;
  \item si $S$ et $S'$ sont deux serveurs adjacents sur l'anneau, toutes les clés de l'intervalle 
   $]h(S), h(S')]$ appartiennent à $S'$.
\end{itemize}
\end{frame}


\begin{frame}
\frametitle{Illustration}

Exemple: l'objet $A$ appartient au serveur IP1-2; l'objet $B$ au serveur (??)  

\vfill

\figSlideWithSize{consist-hashing}{7cm}

\vfill

Si un serveur est ajouté ou retiré, une réorganisation \alert{locale} est suffisante.

\end{frame}


\begin{frame}
\frametitle{Quelques détails vraiment utiles}

Comment fonctionne la tolérance aux pannes\,? Comment fonctionne l'équilibrage\,?


\vfill

% A two-columns presentation
\begin{tabular}{l|p{6cm}}
   \parbox{6cm}{


\alert{Panne} $\Rightarrow$ géré par la réplication (surprise!)\,; 
par exemple on copie sur la machine suivante dans l'anneau, etc.


  \vfill
  
\alert{Equilibrage} $\Rightarrow$ un même serveur (physique) est distribué
en plusieurs points (virtuels) sur l'anneau.

\begin{itemize}
  \item plus il y a de points, plus le serveur recevra de données\,;
  \item en cas de panne, réoganisation répartie plus justement.
\end{itemize}
  
   }  % Inclusion of an XML file
  & 
  \parbox{5cm}{\figSlideWithSize{ch-refinement}{5cm} }
\end{tabular}

\end{frame}

\begin{frame}
\frametitle{Mais où est donc la table de hachage\,?}

\alert{Qui nous dit quelle est la liste des serveurs et leur intervalle}\,?

\smallskip

Il nous faut un \alert{routeur}.

\smallskip

\figSlideWithSize{archi-consistenthash}{8cm} 
\smallskip

Autres solutions\,: (i) la table est répliquée sur chaque n{\oe}ud (acceptable
quand on n'ajoute/retire pas de serveurs tout le temps), (ii) le routeur est intégré au client.



\end{frame}

\subsection{Etude de cas\,: Chord}

\begin{frame}[fragile]
\frametitle{Implementation: Chord Ring}


The hash function $h$ takes the URL of a peer and maps
it to the address space $[0 .. 2^m-1]$. Hashing is based on 
arithmetic modulo $2^m$. 

\begin{center}
h : pId $\rightarrow$ $pId\ mod\ 2^m$
\end{center}

$h$ distributes the peers around a \emph{ring}.
We assume no two peers have the same $h(pId)$ ($m$ large enough). 

\begin{tabular}{m{.4\hsize}m{.4\hsize}}
   \parbox{\hsize}{The Chord Ring with $m=3, 2^m=8$.
Each peer with id $n$ is located
at node $n\ mod\ 8$ on the ring. \\
Ex: peer 22 is assigned to node~6.
}  % Inclusion of an XML file
  & 
   \figSlideWithSize{chord}{5cm} % Inclusion of a PDF file
\end{tabular}

%\vspace*{-2cm}


\end{frame}

\begin{frame}[fragile]
\frametitle{Distribute the keys to the peers}

Keys are hashed with function $h$ with range $[0..2^m-1]$\\

Rule: a key k is assigned to the unique peer p such that  
\begin{itemize}
\item h(p) $\leq$ h(k)
\item \alert{and} there is no p' such that h(p) $<$ h(p') $\leq$ h(k) 
\end{itemize}

We say that p is \emph{responsible} for key k.
\begin{tabular}{m{.4\hsize}m{.5\hsize}}
   \parbox{\hsize}{Assignment of item to peer: item 13
   is hashed to $h(13) = 5$ and assigned to the
   peer p with the largest $h(p) \leq 5$.
}  % Inclusion of an XML file
  & 
   \figSlideWithSize{chord2}{5cm} % Inclusion of a PDF file
\end{tabular}
\end{frame}


\begin{frame}[fragile]
\frametitle{Routing tables}

For each peer p, $friends_p$ contains (at most) $\log 2^m = m$ peer addresses.

For each i in [1..m], the i$^{th}$ friend p$_i$ is such that 
\begin{itemize}
\item h(p$_i$) $\leq$ h(p)+$2^{i-1}$ 
\item there is no p' such that h(p$_i$) $<$ h(p') $\leq$
h(p)+$2^{i-1}$ 
\end{itemize}

In other words:  p$_i$ is the peer responsible  for key $h(p) + 2^{i-1}$.

\begin{examples}[Let $m=10$, $2^m = 1024$; consider peer p with $h(p) =
10$.]

\begin{itemize}
  \item The first friend p$_1$ is the peer resp.  for $10 + 2^0 = 11$
 \item The second friend p$_2$ is the peer resp.  for  $10 + 2^1 = 12$
 \item The third friend p$_3$ is the peer resp.  for  $10 + 2^2 = 14$
 \item  friend p$_7$ is the peer resp.  for  $10 + 2^6 = 74$
 \item  The last friend p$_{10}$ is the peer resp.  for  $(10 + 512) = 522$
\end{itemize}
\end{examples}
Exercise: express the gap between friends[i] and friends[i+1]

\end{frame}

\begin{frame}
\frametitle{Understanding routing tables}

Important properties:
\begin{enumerate}
  \item a peer maintains a small routing table, e.g., 16 friends
          for each peer, in a ring with $2^{16} = 65,536$ nodes;
  \item each peer knows better the peers close on the ring that the peers
        far away;
   \item a peer p cannot (in general) find directly the peer p'
      responsible for a key k; \emph{but p can find a friend which holds a more
      accurate information about k.}
\end{enumerate}

\begin{tabular}{m{.4\hsize}m{.4\hsize}}
   \parbox{\hsize}{The figure shows the friends of
   peer p16, located at node 0 (note the collisions). p16
   does \emph{not} know p22.\\
   
   Exercise: what are the friends of p11, located at node 3? 
}  % Inclusion of an XML file
  & 
   \figSlideWithSize{chord-friends}{5cm} % Inclusion of a
\end{tabular}

\end{frame}

\begin{frame}[fragile]
\frametitle{Searching Chord: The \textit{get} algorithm}

Peer p looks for index(k):
\begin{itemize}
  \item if p is responsible for k, we are done;
\item else  let $i$ such that  h(p)+$2^{i-1}$ $\leq$ h(k) < h(p)+$2^{i}$ :
forward the search to the friend p$_i$.
\end{itemize}

Fact: the search range is (at worse) halved at each step $\rightarrow$ the
search converges in $O(\log 2^m) = O(m)$ hops.


\begin{examples}[p16 receives a request for item k, with
      h(k)=7.]

\begin{tabular}{m{.5\hsize}m{.4\hsize}}
   \begin{enumerate}
     \item p16 forwards the request to p11, its third friend (why?).
     \item then p11 forwards to p22, its third friend (why?).
     \item p22  find item k locally.
   \end{enumerate}
      
Exercise:  develop an example of the worst case with $m=10$.
 % Inclusion of an XML file
  & 
   \figSlideWithSize{chord-get}{5cm} % Inclusion of a
\end{tabular}
\end{examples}
\end{frame}


\end{document}


